
\documentclass[11pt]{article}
\usepackage[a4paper,margin=1in]{geometry}
\usepackage{amsmath,amssymb}
\usepackage{lmodern}
\usepackage[T1]{fontenc}
\usepackage{microtype}
\usepackage{hyperref}
\setlength{\parindent}{0pt}
\setlength{\parskip}{0.65em}
\allowdisplaybreaks

\begin{document}

Let \(N\) be the population size. Let \(\tau\) be dimensionless time measured in units of the mean infectious period. Let \(x(\tau)\) be the susceptible fraction, \(y_1(\tau)\) the infected fraction of strain 1, and \(y_2(\tau)\) the infected fraction of strain 2. Let \(\varepsilon\) be the demographic timescale parameter. Let \(\mathcal R_{0,1}\) and \(\mathcal R_{0,2}\) be the basic reproduction numbers of strain 1 and strain 2.

The deterministic two-strain system is
\[
\frac{dx}{d\tau}
=
\varepsilon(1-x)-\mathcal R_{0,1}xy_1-\mathcal R_{0,2}xy_2,
\]
\[
\frac{dy_1}{d\tau}
=
(\mathcal R_{0,1}x-1)y_1,
\]
\[
\frac{dy_2}{d\tau}
=
(\mathcal R_{0,2}x-1)y_2.
\]
where
\[
\mathcal R_{0,1}>1,
\qquad
\mathcal R_{0,2}>1.
\]

The strain-specific boundary-layer thresholds are defined by
\[
y_1^\star
=
\varepsilon\left(1-\frac{1}{\mathcal R_{0,1}}\right),
\]
\[
y_2^\star
=
\varepsilon\left(1-\frac{1}{\mathcal R_{0,2}}\right).
\]

Assume strain 1 enters its boundary layer before strain 2. The opposite case is obtained by interchanging the labels 1 and 2.

Let \(\tau_1\) be the time at which strain 1 enters its boundary layer on the declining side of the epidemic wave. Define
\[
x_a=x(\tau_1),
\]
\[
y_1(\tau_1)=y_1^\star,
\]
\[
y_{2,a}=y_2(\tau_1).
\]
By assumption, strain 2 has not yet entered its own boundary layer at this time, so
\[
y_{2,a}>y_2^\star.
\]
The number of strain-1 infectives at this boundary-layer entrance is
\[
m_1=Ny_1^\star.
\]

Define local time after strain 1 enters its boundary layer by
\[
s=\tau-\tau_1.
\]
Thus, at \(s=0\),
\[
x(0)=x_a,
\]
\[
y_2(0)=y_{2,a}.
\]

During the first stage, strain 1 is rare while strain 2 remains in the deterministic regime. We approximate strain 1 by a branching process and strain 2 by a deterministic trajectory. In this stage we ignore the feedback of strain 1 on the susceptible pool, and compared with the strain-2 transmission term, $\varepsilon(1-x)$ can be ignored during this phase. Thus the resident strain-2 trajectory is approximated by
\[
\frac{dx}{ds}
=
-\mathcal R_{0,2}xy_2,
\]
\[
\frac{dy_2}{ds}
=
(\mathcal R_{0,2}x-1)y_2.
\]

By the chain rule,
\[
\frac{dy_2}{dx}
=
\frac{dy_2/ds}{dx/ds}.
\]
Substituting the two equations above gives
\[
\frac{dy_2}{dx}
=
\frac{(\mathcal R_{0,2}x-1)y_2}{-\mathcal R_{0,2}xy_2}.
\]
Cancelling \(y_2\),
\[
\frac{dy_2}{dx}
=
-1+\frac{1}{\mathcal R_{0,2}x}.
\]
Integrating from \((x_a,y_{2,a})\) to \((x,y_2(x))\), using \(u\) as the integration variable,
\[
y_2(x)-y_{2,a}
=
\int_{x_a}^{x}
\left(-1+\frac{1}{\mathcal R_{0,2}u}\right)du.
\]
Therefore
\[
y_2(x)
=
y_{2,a}+x_a-x+\frac{1}{\mathcal R_{0,2}}\ln\frac{x}{x_a}.
\]

Let \(x_b\) be the susceptible fraction at which strain 2 enters its boundary layer. It is determined by
\[
y_2(x_b)=y_2^\star.
\]
Substituting the expression for \(y_2(x)\),
\[
y_{2,a}+x_a-x_b+\frac{1}{\mathcal R_{0,2}}\ln\frac{x_b}{x_a}
=
y_2^\star.
\]
Define
\[
C_2=y_{2,a}+x_a-y_2^\star.
\]
Then \(x_b\) satisfies
\[
x_b-\frac{1}{\mathcal R_{0,2}}\ln\frac{x_b}{x_a}
=
C_2.
\]
Equivalently, using the principal branch \(W_0\) of the Lambert \(W\) function,
\[
x_b
=
-\frac{1}{\mathcal R_{0,2}}
W_0
\left(
-\mathcal R_{0,2}x_a\exp(-\mathcal R_{0,2}C_2)
\right).
\]
At this point, strain 2 has
\[
m_2=Ny_2^\star
\]
infectives.

We now describe the distribution of strain 1 at the time strain 2 enters its boundary layer. Let \(Z_1(s)\) be the number of strain-1 infectives during the first stage. Since strain 1 is rare, we approximate \(Z_1(s)\) by a time-inhomogeneous linear birth-death process. At local time \(s\), the per capita birth rate of strain 1 is
\[
b_1(s)=\mathcal R_{0,1}x(s),
\]
and the per capita death rate is
\[
d_1(s)=1.
\]

For a process starting from one strain-1 infective, define the probability generating function
\[
G_1(s;z)
=
\mathbb E[z^{Z_1(s)}\mid Z_1(0)=1],
\]
where \(z\) is the generating-function variable.

For a time-inhomogeneous linear birth-death process with per capita birth rate \(b(s)\) and death rate \(d(s)\), the probability generating function satisfies the backward equation
\[
\frac{\partial G}{\partial s}
=
b(s)G^2-[b(s)+d(s)]G+d(s),
\]
with initial condition
\[
G(0;z)=z.
\]
In our case,
\[
\frac{\partial G_1}{\partial s}
=
\mathcal R_{0,1}x(s)G_1^2
-
[\mathcal R_{0,1}x(s)+1]G_1
+
1,
\]
with
\[
G_1(0;z)=z.
\]

This is the Riccati equation. Define the cumulative net growth
\[
H_1(s)
=
\int_0^s[\mathcal R_{0,1}x(u)-1]du,
\]
where \(u\) is an integration variable. Define 
\[
A_1(s)
=
\int_0^s \mathcal R_{0,1}x(v)\exp(H_1(v))dv,
\]
where \(v\) is an integration variable. Then the solution of the Riccati equation is
\[
G_1(s;z)
=
1-
\frac{\exp(H_1(s))(1-z)}
{1+(1-z)A_1(s)}.
\]

We now rewrite these quantities in terms of \(x\), avoiding the need to compute the time at which strain 2 enters its boundary layer. From
\[
\frac{dx}{ds}
=
-\mathcal R_{0,2}xy_2(x),
\]
we have
\[
ds
=
-\frac{dx}{\mathcal R_{0,2}x y_2(x)}.
\]
When the system moves from \(x_a\) to a smaller value \(u\), define
\[
\widehat H_1(u)
=
\int_u^{x_a}
\frac{\mathcal R_{0,1}v-1}
{\mathcal R_{0,2}v\,y_2(v)}
\,dv,
\]
where \(v\) is an integration variable. In particular,
\[
\widehat H_1(x_b)
=
\int_{x_b}^{x_a}
\frac{\mathcal R_{0,1}v-1}
{\mathcal R_{0,2}v\,y_2(v)}
\,dv.
\]
And define
\[
\widehat A_1
=
\int_{x_b}^{x_a}
\frac{\mathcal R_{0,1}}
{\mathcal R_{0,2}y_2(u)}
\exp(\widehat H_1(u))
\,du,
\]
where \(u\) is an integration variable.

Therefore, from one strain-1 infective at \(x=x_a\), the probability generating function for the number of strain-1 infectives when strain 2 enters its boundary layer is
\[
G_1^{1\mid 2}(z)
=
1-
\frac{\exp(\widehat H_1(x_b))(1-z)}
{1+(1-z)\widehat A_1}.
\]
The superscript \(1\mid 2\) indicates that strain 1 is rare while strain 2 is resident.

If there are \(m_1\) strain-1 infectives when strain 1 first enters its boundary layer, the branching property gives the total generating function
\[
\left[G_1^{1\mid 2}(z)\right]^{m_1}.
\]

Now consider the second stage where both strains are in their boundary layers. Let \(r\) be local time starting from the moment strain 2 enters its boundary layer. At \(r=0\),
\[
x(0)=x_b.
\]
In this common boundary layer, both strains are rare. We ignore the feedback of infections on the susceptible pool, so susceptible recovery is dominated by demographic replenishment only:
\[
\frac{dx}{dr}
=
\varepsilon(1-x).
\]
Solving this equation gives
\[
x(r)
=
1-(1-x_b)\exp(-\varepsilon r).
\]

In the common boundary layer, the two strains are approximated as conditionally independent birth-death processes in the same externally determined susceptible environment \(x(r)\). For strain \(i\), where \(i=1,2\), the per capita birth rate is
\[
b_i^C(r)=\mathcal R_{0,i}x(r),
\]
and the per capita death rate is
\[
d_i^C(r)=1.
\]
The superscript \(C\) denotes the common boundary layer.

Define the cumulative net growth for strain \(i\) in the common boundary layer by
\[
H_i^C(r)
=
\int_0^r[\mathcal R_{0,i}x(u)-1]du.
\]
Since
\[
x(u)=1-(1-x_b)\exp(-\varepsilon u),
\]
we obtain
\[
H_i^C(r)
=
(\mathcal R_{0,i}-1)r
-
\frac{\mathcal R_{0,i}(1-x_b)}{\varepsilon}
(1-\exp(-\varepsilon r)).
\]
And define 
\[
A_i^C(r)
=
\int_0^r
\mathcal R_{0,i}x(v)\exp(H_i^C(v))dv.
\]
The probability generating function at time \(r\) is
\[
G_i^C(r;z)
=
1-
\frac{\exp(H_i^C(r))(1-z)}
{1+(1-z)A_i^C(r)}.
\]
Therefore the eventual extinction probability for one strain-\(i\) infective in the common boundary layer is
\[
q_i^C
=
\lim_{r\to\infty}G_i^C(r;0).
\]
Equivalently,
\[
q_i^C
=
1-
\lim_{r\to\infty}
\frac{\exp(H_i^C(r))}
{1+A_i^C(r)}.
\]

Then we can get the burnout probability for strain 2. Strain 2 enters its boundary layer with
\[
m_2=Ny_2^\star
\]
infectives. Each infective lineage has eventual extinction probability \(q_2^C\). Hence
\[
Q_2
=
(q_2^C)^{m_2}.
\]

For strain-1, if there are \(n\) strain-1 infectives at that moment, the probability that all of them eventually go extinct in the common boundary layer is
\[
(q_1^C)^n.
\]
Here \(n\) is the random variable \(Z_1(x_b)\).  We take the expectation
\[
\mathbb E[(q_1^C)^{Z_1(x_b)}].
\]
By the definition of a probability generating function, this is
\[
G_1^{1\mid 2}(q_1^C)
\]
for a single initial strain-1 infective at \(x=x_a\). Since strain 1 enters its boundary layer with \(m_1\) infectives, the strain-1 burnout probability for the whole wave is
\[
Q_1
=
\left[G_1^{1\mid 2}(q_1^C)\right]^{m_1}.
\]

Thus the one-wave burnout probabilities under the two-boundary approximation are
\[
Q_1
=
\left[G_1^{1\mid 2}(q_1^C)\right]^{m_1},
\]
\[
Q_2
=
(q_2^C)^{m_2}.
\]


\end{document}
